% Source-preserving integrated proof body; input from toe_round7_proofs.tex.
\subsubsection*{Round 7: first-order matter coupling of the staggered tensor target}
\label{sec:toe-r7-matter}
\noindent\emph{The displayed arguments below are the proof-bearing source.  The current
repository regression is \veri{v1035\_\allowbreak matter\_\allowbreak coupling.py}; its finite checks audit the
implementation and do not replace the universal steps in the proof.}

\begingroup
\def\TT{\mathrm{TT}}
\def\Tr{\operatorname{Tr}}
\def\ii{\mathrm{i}}
\def\dd{\mathrm{d}}
\def\cO{\mathcal O}
\def\cJ{\mathcal J}
\begin{center}
\fbox{\parbox{0.93\linewidth}{
\textbf{Result and boundary.}  Conditional on the declared staggered Fourier
dictionary, every nonzero momentum admits an exact affine coupling to a
prescribed classical source obeying the two Ward equations.  Source energy and
momentum must deform the scalar and vector constraints; with those terms, the
fibre constraints propagate and the curvature has the classical retarded
response displayed below.  The executable certificate does not import
repository code, but it does instantiate the declared tensor/vector masks and
checks their three reciprocal-axis Bloch transitions.  Dynamical
matter is tested only in collinear one-dimensional sectors and only through
first order in the coupling.  A free scalar and a fermion with a
nearest-neighbour one-particle Hamiltonian pass those tests after their
composite lattice vertices are repaired.  A naive scalar stress fails, and an
onsite $\phi^4$ interaction has an explicit homogeneous momentum defect that no
periodic stress divergence can remove.  The deleted homogeneous gravity block
cannot absorb nonzero total matter energy or momentum.  No full three-dimensional
matter coupling, quantum Kubo calculation, nonlinear constraint algebra,
interacting quantum gravity, Lorentz-covariant lattice Riemann tensor, TFPT
parent, T7 closure, TOE claim, or RH claim is proved.
}}
\end{center}

\paragraph{Fixed regulator and conventions.}

We state the spatial regulator of the preceding classical and quantum
artifacts as the intended real-space interpretation.  The checker below
instantiates these declared masks algebraically, but does not import a
repository implementation.  On a periodic cubic lattice of spacing $a$, a component with mask
$m\in\{0,1\}^3$ lives at $x+am/2$, and
\[
 D_i^{(m)}f(x)=\frac1a
 \begin{cases}
 f(x+ae_i)-f(x),&m_i=0,\\
 f(x)-f(x-ae_i),&m_i=1.
 \end{cases}                                      \tag{1}
\]
Thus $D_i^{(m\mathbin{\rm xor}e_i)}=-(D_i^{(m)})^T$ under the periodic
lattice sum.  Physical-position Fourier phases give symbol
$\ii\kappa_i$; the factor $\ii$ is retained below, where
\[
 \kappa_i=\frac2a\sin\frac{ak_i}{2},\qquad
 r^2=\kappa^T\kappa.                              \tag{2}
\]
The diagonal tensor components have mask zero, $q_{ij}$ with $i\ne j$ has
mask $e_i\mathbin{\rm xor}e_j$, and a vector component $j_i$ has mask $e_i$.
Products below are placed on the indicated cells or links; no unmentioned
continuum Leibniz rule is used.

In the six orthonormal symmetric-tensor coordinates
$(11,22,33,23,13,12)$ define the following real matrices; they do not include
the Fourier factor $\ii$:
\[
 (Bh)_i=h_{ij}\kappa_j,\quad
 t=(1,1,1,0,0,0)^T,\quad v=B^T\kappa,\quad
 s=v-r^2t,                                         \tag{3}
\]
\[
 K_p=I-\frac12tt^T,\qquad
 K_q=r^2(I-tt^T)-2B^TB+tv^T+vt^T.                 \tag{4}
\]
With the common Fourier-volume normalization suppressed, the real canonical
Hamiltonian sums the paired complex fibres; the vacuum constraints at each
nonzero momentum are
\[
 \begin{gathered}
 H_0=\frac12\sum_{k\ne0}
 \bigl[p(-k)^TK_p(k)p(k)+q(-k)^TK_q(k)q(k)\bigr],\\
 c_0=-s^Tq,\qquad \boldsymbol c=\ii Bp.
 \end{gathered}                                      \tag{5}
\]
Relative to the Round-6 real matrices, these are only nonzero sign/phase
changes of the vacuum constraints.  Their zero surface, rank, and the reduced
$E,F$ observables are therefore unchanged.
The inherited exact identities are
\[
 K_qB^T=0,\quad K_ps=B^T\kappa,\quad K_qt=s,
\quad K_qK_pK_q=r^2K_q+\frac12ss^T=r^4P_{\TT}.    \tag{6}
\]
Here $c_0$ is the Fourier transform of
$D_iD_jq_{ij}-\Delta\operatorname{tr}q$, and $c_i$ is the transform of
$D_jp_{ij}$.  Thus
$\dot c_0=\ii\kappa_i c_i=-\kappa^TBp$ and
$\dot {\boldsymbol c}=0$.  Dropping $\ii$ and restoring it only after the algebra
gives the wrong scalar sign, as the plane-wave mutant below demonstrates.

\paragraph{The coupling that actually preserves the constraints.}

Let $(\rho,j_i,\tau_{ij})$ be a symmetric c-number source on the scalar,
vector, and tensor staggered cells.  Its two Ward equations are
\[
 \dot\rho+D_i j_i=0,\qquad
 \dot j_i+D_j\tau_{ij}=0.                          \tag{7}
\]
At $k\ne0$, add the spatial metric coupling and deform both constraints:
\[
 H_g(t)=H_0+g\langle q,\tau(t)\rangle,\qquad
 \Psi_0=-s^Tq+g\rho,\qquad
 \boldsymbol\Psi=\ii Bp-g\boldsymbol j.                          \tag{8}
\]
The lapse and shift multipliers multiply $\Psi_0$ and $\boldsymbol\Psi$ in the total
Hamiltonian.  Thus (8) is not merely the isolated $q_{ij}\tau_{ij}$ vertex:
the $\rho$ and $j_i$ couplings are present as constraint terms.  Periodic
summation by parts fixes all signs in (8).

\begin{theorem}[Conditional affine source consistency in the complex Fourier fibre]
Assume that (3)--(6) are the fibre dictionary of the intended staggered
placement and that a prescribed c-number source satisfies (7).  Then the four
constraints in (8) are first class within this fibre algebra and propagate as
\[
 \dot\Psi_0=D_i\Psi_i,\qquad \dot\Psi_i=0.        \tag{9}
\]
If one instead retains the vacuum vector constraint, its time derivative is
\[
 \dot c_i=-gD_j\tau_{ij}=g\dot j_i,                \tag{10}
\]
which is generically nonzero.  Consequently ``couple $q_{ij}$ to stress''
without deforming the constraints is inconsistent already at order $g$.
Moreover, within the linear ansatz
\[
 H_{\rm int}=\gamma g\langle q,\tau\rangle,
 \quad\Psi_0=c_0+\alpha g\rho,
 \quad\boldsymbol\Psi=\boldsymbol c+\beta g\boldsymbol j,
\]
Ward propagation uniquely fixes $\alpha=\gamma$ and $\beta=-\gamma$.
Thus (8), up to the overall normalization of $g$, is the smallest affine
constraint repair of the spatial stress vertex.
\end{theorem}

\begin{proof}
Hamilton's equations give
$\dot q=K_pp$ and $\dot p=-K_qq-g\tau$.  At fixed momentum, (7) reads
$\dot\rho=-\ii\kappa^Tj$ and $\dot j=-\ii B\tau$.  Therefore
\[
 \dot\Psi_0=-\kappa^TBp-\ii g\kappa^Tj
             =\ii\kappa^T(\ii Bp-gj)=\ii\kappa^T\boldsymbol\Psi,
\]
\[
 \dot{\boldsymbol\Psi}=-\ii gB\tau-g\dot j=0.
\]
This is (9) with no stripped-$\ii$ intermediate convention.  Since the sources are c-numbers, their Poisson
brackets vanish and the affine shifts do not change the vacuum first-class
brackets.  Dropping $-g\boldsymbol j$ leaves (10).  For the general coefficients,
the vector residual is $-\ii(\gamma+\beta)gB\tau$, while the scalar residual
relative to $\ii\kappa^T\boldsymbol\Psi$ is
$-\ii(\alpha+\beta)g\kappa^Tj$; both vanish only for the stated coefficients.
\end{proof}

The sign is falsifiable on one mode.  Take the plane wave
$e^{-\ii t+\ii z}$ with $\kappa=e_3$ and amplitudes
$\rho=j_3=\tau_{33}=g=1$.  Both Ward equations hold.  Choose
$Bp=-\ii e_3$ and $s^Tq=1$, so the repaired constraints vanish.  Directly,
$\dot\Psi_0=\dot{\boldsymbol\Psi}=0$.  If one instead keeps the old stripped-$\ii$
choice $\Psi_0=+s^Tq+g\rho$ while restoring
$\boldsymbol\Psi=\ii Bp-gj$, the propagation residual is
$\dot\Psi_0-\ii\kappa^T\boldsymbol\Psi=-2\ii$.

The source need not be positive, local in time, or generated by matter for
this theorem.  Fourier reality is explicit:
$B(-\kappa)=-B(\kappa)$ while $s,K_p,K_q$ are even.  Hence
$\Psi_0(-k)=\overline{\Psi_0(k)}$ and
$\boldsymbol\Psi(-k)=\overline{\boldsymbol\Psi(k)}$ whenever the fields at $-k$ are the
conjugates of those at $k$ and $g$ is real.  For an unordered non-self-conjugate
pair, the actual real Hamiltonian contribution is
\[
 H_{\{k,-k\}}=p(-k)^TK_p\,p(k)+q(-k)^TK_q\,q(k)
 +g\bigl[q(-k)^T\tau(k)+q(k)^T\tau(-k)\bigr],
\]
up to the common Fourier-volume normalization.  It equals its complex
conjugate.  To lift the result from the
complex fibre calculation to the actual staggered lattice it must also obey
the mask placements and Bloch gluing.

That gluing is checked explicitly for the declared regulator.  Under the
reciprocal shift $k\mapsto k+2\pi e_a/a$, let $R_V^{(a)}$ change the sign of
vector component $a$, and let $G_T^{(a)}$ change the sign of an off-diagonal
tensor component precisely when its mask contains $e_a$.  Diagonal tensor
components are unchanged.  Then
\[
 \kappa'=R_V^{(a)}\kappa,\quad
 B'G_T^{(a)}=R_V^{(a)}B,\quad s'=G_T^{(a)}s,
 \quad K_q'=G_T^{(a)}K_qG_T^{(a)}.
\]
Since $K_p$ commutes with $G_T^{(a)}$, the source gluing
$\tau'=G_T^{(a)}\tau$, $j'=R_V^{(a)}j$, and $\rho'=\rho$ gives
\[
 \cJ'=G_T^{(a)}\cJ,\qquad
 \Psi_0'=\Psi_0,\qquad \boldsymbol\Psi'=R_V^{(a)}\boldsymbol\Psi,
 \qquad (q')^T\tau'=q^T\tau .
\]
Here $\cJ$ is the response drive defined in (12).  The checker verifies all
these matrix identities for $a=1,2,3$.  Reality is
most cleanly imposed on extended-zone labels $(-k,k)$; if $-k$ is folded into
the first Brillouin zone by a reciprocal vector, the displayed transition
matrices supply the additional mask signs.  What remains unverified is that a
particular repository implementation uses exactly these declared conventions.

\paragraph{Gauge-invariant retarded curvature response.}

The local curvature variables remain
\[
 E=K_qq,\qquad F=K_qK_pp.                         \tag{11}
\]
They commute strongly with the gravitational parts of all four constraints:
$K_qB^T=0$ and $K_qK_ps=0$.  The c-number shifts in (8) have zero bracket
with them, so $E,F$ remain gauge invariant.  They obey $\dot E=F$ and, on
$\Psi_0=0$,
\[
 \ddot E+r^2E=-g\cJ,\qquad
 \cJ:=K_qK_p\tau+\frac12s\rho.                  \tag{12}
\]
The retarded solution with vanishing past data is therefore
\[
 \boxed{\quad
 \delta E(t,k)=-g\int_{-\infty}^{t}
 \frac{\sin(r(t-t'))}{r}\,\cJ(t',k)\,\dd t',
 \qquad \delta F=\partial_t\delta E .\quad}       \tag{13}
\]
This is the classical Hamiltonian source response in the complex Fourier
dictionary, not a spectral-positivity surrogate.  Identifying it with a quantum
Kubo response requires importing the Round-6 CCR/Fock representation and
evaluating the corresponding commutator; that calculation is not executed in
this artifact.

\begin{theorem}[Proposition: Ward compatibility of the response]
The source in (12) is transverse and has the trace required by the sourced
scalar constraint:
\[
 B\cJ=0,\qquad t^T\cJ=-(\ddot\rho+r^2\rho).       \tag{14}
\]
Consequently (13), with constraint-compatible initial data, satisfies
\[
 BE=0,\qquad t^TE=g\rho,\qquad t^TF=g\dot\rho.  \tag{15}
\]
For a TT step source $\tau(t)=\Theta(t)e_+$ at
$\kappa=(0,0,r)$ and $\rho=j=0$,
\[
 \delta E(t)=-g(1-\cos rt)e_+,\qquad t\ge0,       \tag{16}
\]
where $e_+=(1,-1,0,0,0,0)^T/\sqrt2$.
\end{theorem}

\begin{proof}
Both $BK_q=0$ and $Bs=BK_qt=0$, proving transversality.  From (6),
\[
 t^TK_qK_p\tau=s^TK_p\tau=\kappa^TB\tau=-\ddot\rho,
 \qquad t^Ts=-2r^2,
\]
where the second Ward equation followed by the time derivative of the first
was used.  This proves (14).  Taking the trace of (12) shows that
$t^TE-g\rho$ obeys the homogeneous wave equation and hence vanishes for
compatible initial data.  Equation (16) follows from $K_qK_pe_+=r^2e_+$.
\end{proof}

Notice that sourced $E$ is not generally TT: it is transverse, while its trace
is fixed by the energy density.  Only its radiative part is TT.  Calling all of
$E$ a free graviton after matter is added would erase the scalar constraint.

\paragraph{Dynamical matter: what is and is not proved.}

Let matter have canonical variables $z$ and an uncoupled Hamiltonian $H_m$.
If its zeroth-order densities obey (7), then
\[
 H=H_0+H_m+g\langle q,\tau^{(0)}(z)\rangle,\qquad
 \Psi_0=-s^Tq+g\rho^{(0)}(z),\quad
 \boldsymbol\Psi=\ii Bp-g\boldsymbol j^{(0)}(z)                 \tag{17}
\]
obeys (9) modulo $\cO(g^2)$.  This is the precise first-order feedback result.
It is not exact nonlinear closure.  Indeed, even before computing the
$g$-correction to the matter stress,
\[
 \{\Psi_0(x),\Psi_i(y)\}
 =-g^2\{\rho^{(0)}(x),j_i^{(0)}(y)\}_m,           \tag{18}
\]
plus analogous current brackets.  These terms are generally nonzero and must
be absorbed by second-order constraint and Hamiltonian corrections if a
nonlinear theory exists.  The checker gives a canonical scalar datum for which
the matter bracket on the right of (18) equals $1/2$ exactly.

The distinction between a canonical/Noether stress and the metric-response
stress matters here, especially for spinors.  Their equivalence requires
hypotheses and, for spinors, a tetrad formulation; see Leclerc's primary
analysis \cite{toe-r7-matter-Leclerc}.  We therefore test only explicit lattice Ward
complexes below and do not rename them a full Hilbert stress tensor.

\paragraph{Free scalar: a local repair, then a sharp $\phi^4$ obstruction.}

It suffices to examine a collinear sector; failure there is failure of a
purported universal cubic-lattice identity.  Set $a=1$ on a periodic
one-dimensional row and write
\[
 g_x=\phi_{x+1}-\phi_x,\qquad
 \Delta\phi_x=g_x-g_{x-1},\qquad
 \dot\phi_x=\pi_x,\quad
 \dot\pi_x=\Delta\phi_x-m^2\phi_x.                \tag{19}
\]
The site energy and link flux
\[
 \rho_x=\frac12\pi_x^2+\frac12m^2\phi_x^2
 +\frac14(g_x^2+g_{x-1}^2),\qquad
 j_x=-\frac12(\pi_x+\pi_{x+1})g_x                \tag{20}
\]
obey the exact energy equation
\[
 \dot\rho_x+j_x-j_{x-1}=0.                        \tag{21}
\]
The continuum-looking averaged-gradient stress
\[
 \tau_x^{\rm naive}=\frac12\pi_x^2
 +\frac14(g_x^2+g_{x-1}^2)-\frac12m^2\phi_x^2    \tag{22}
\]
does not obey momentum continuity.  An exact same-range correction is
\[
 \boxed{\quad
 \tau_x^{\rm free}=\frac12\pi_x^2+\frac12g_xg_{x-1}
 -\frac12m^2\phi_x^2
 =\tau_x^{\rm naive}-\frac14(\Delta\phi_x)^2 .\quad} \tag{23}
\]
Direct substitution gives
\[
 \dot j_x+\tau_{x+1}^{\rm free}-\tau_x^{\rm free}=0. \tag{24}
\]
The correction in (23) is exactly the finite-difference product-rule term;
it vanishes at leading continuum order but is mandatory at finite spacing.

For the rational four-site datum
$(\phi_0,\phi_1,\phi_2,\phi_3)=(0,1,2,4)$, the four residuals obtained from
(22) are
\[
 \left(-\frac{25}{4},\frac14,\frac{35}{4},-\frac{11}{4}\right),             \tag{25}
\]
so this is an analytic mutant, not a rounding effect.

Now add $V_{\rm int}=\lambda\phi^4/24$ to $\rho_x$, the force
$-\lambda\phi_x^3/6$, and the continuum-like pressure
$-\lambda\phi_x^4/24$ to (23).  Energy continuity (21) remains exact, but
the momentum residual is
\[
 R_x:=\dot j_x+\tau_{x+1}-\tau_x
 =\frac{\lambda}{24}(\phi_x+\phi_{x+1})
   (\phi_{x+1}-\phi_x)^3.                          \tag{26}
\]
For the same datum,
\[
 \sum_{x=0}^{3}R_x=-\frac{17}{2}\lambda\ne0.      \tag{27}
\]
Every periodic stress divergence sums to zero.  Therefore no redefinition of
$\tau_x$ alone---local or nonlocal---can repair (26) on this matter phase
space while leaving $j_x$ and its dynamics fixed.  This is the first precise
interacting obstruction, and it appears in constraint propagation at order
$g\lambda$, not only in a higher-order gravity bracket.

There is a limited repair for the already-selected nonzero gravity modes.  Let
$P_0$ be the spatial mean and let $D\tau=\tau_{x+1}-\tau_x$.  Then
\[
 \delta\tau=-D^+(I-P_0)R,\qquad
 D\delta\tau=-(I-P_0)R,                            \tag{28}
\]
where $D^+$ is any zero-mean inverse of $D$.  This makes the nonzero-mode Ward
identity exact and is constructed explicitly by cumulative summation in the
checker.  It is nonlocal, leaves the homogeneous momentum defect $P_0R$, and
does not constitute a local interacting matter coupling.  A full repair must
instead alter the matter discretization/current or add degrees of freedom that
carry the missing lattice momentum.  The broader lattice literature likewise
defines and renormalizes energy-momentum tensors by enforcing Ward identities,
rather than assuming the continuum expression survives unchanged
\cite{toe-r7-matter-BuonannoCellaCurci,toe-r7-matter-GiustiPepe}.

\paragraph{Free fermion: exact collinear divided-difference vertices.}

A second collinear test uses a free chiral fermion whose \emph{one-particle
Hamiltonian} is nearest-neighbour, with dispersion $\epsilon(p)=\sin p$.
This is a finite-spacing Ward test, not a claim of a doubler-free $3+1$
dimensional chiral theory.  For a
bilinear transferring momentum $K$, set
\[
 \kappa(K)=2\sin\frac K2,\quad
 \epsilon_- =\sin p,\quad \epsilon_+=\sin(p+K),\quad
 \bar\epsilon=\frac12(\epsilon_++\epsilon_-).     \tag{29}
\]
The exact divided-difference identity is
\[
 \epsilon_+-\epsilon_-
 =\kappa(K)\cos\left(p+\frac K2\right).            \tag{30}
\]
Consequently the energy, current, and stress vertices
\[
 \rho=\bar\epsilon,\quad
 j=\cos\left(p+\frac K2\right)\bar\epsilon,\quad
 \tau=\cos^2\left(p+\frac K2\right)\bar\epsilon  \tag{31}
\]
satisfy both one-particle Ward identities
\[
 (\epsilon_+-\epsilon_-)\rho=\kappa j,\qquad
 (\epsilon_+-\epsilon_-)j=\kappa\tau.             \tag{32}
\]
In real space, $\rho$, $j$, and $\tau$ in (31) have maximal hopping ranges
one, two, and three, respectively: only the one-particle Hamiltonian is
nearest-neighbour.  The average factors are the fermionic analogue of the
scalar repair in (23).
Replacing $j$ by the continuum vertex $\rho$ fails: at
$p=\pi/6$, $K=\pi/3$, one has
$(\rho,j,\tau)=(3/4,3/8,3/16)$ and the first Ward residual is $-3/8$.
Thus both bosonic and fermionic tests reject naive continuum vertices, while
both free tests admit explicit finite-spacing corrections.

\paragraph{The homogeneous obstruction.}

The preceding free tensor construction removed the complete $k=0$ canonical
block to eliminate its negative trace direction.  At $k=0$,
$B=s=K_q=0$, so its vacuum local constraints vanish identically.  Equation
(8) would reduce to
\[
 \Psi_0(0)=g\rho(0),\qquad \boldsymbol\Psi(0)=-g\boldsymbol j(0). \tag{33}
\]
A positive-energy matter state normally has $\rho(0)\ne0$, and a state may
also have $\boldsymbol j(0)\ne0$.  For $g\ne0$, imposing (33) would delete either
such component rather than source gravity.  The conditional coupling must
therefore be read in one of two narrow ways:
\begin{enumerate}[label=(\alph*),leftmargin=2em]
\item couple only $(I-P_0)(\rho,j,\tau)$, equivalently fluctuations around a
specified homogeneous background; or
\item restore and consistently stabilize a homogeneous/global geometry sector.
\end{enumerate}
The first option is sufficient for (9)--(16) at nonzero momentum but does not
describe the gravitational response to total energy.  The second is an open
construction.  Simply calling the removed block ``gauge'' is not a repair.

\paragraph{Quantum meaning and remaining obligation.}

Equation (13) is an exact classical Green-function solution of the conditional
staggered Fourier equations and is gauge invariant and Ward compatible there.  A
smooth compact-time c-number source is expected to define linear forcing of
the already-constructed oscillator Fock space, but the CCR import and Kubo
commutator needed to prove that quantum statement are not executed here.  Thus
the generic prescribed-source fibre algebra is closed; its full staggered
real-space and quantum realization remain open checks.

For dynamical scalar or fermion matter, (17) is an interacting Hamiltonian on
the tensor--matter tensor product only as a formal first-order deformation.
Equations (18), (26), and (33) are three distinct unclosed edges:
\begin{itemize}[leftmargin=2em]
\item $g^2$ current brackets require genuinely deformed constraints and
gravity self-couplings;
\item finite-spacing interacting matter may fail the spatial Ward identity
already at order $g\lambda$;
\item total energy has no receiving homogeneous gravity mode.
\end{itemize}
An exact graph energy current such as $J_{x\to y}=\ii[h_x,h_y]$ addresses
energy continuity only.  It is not by itself a symmetric stress satisfying
the vector Ward identity, and it does not close the Hamiltonian constraint
algebra.  Conversely, positivity of a TT stress spectral measure says nothing
about these three equations.

\paragraph{Executable certificate and evidence typing.}

The repository certificate \veri{v1035\_\allowbreak matter\_\allowbreak coupling.py} imports no proof text.
It checks the generic prescribed-source polynomial identities with their full
complex Fourier signs, a conserved plane-wave witness, the mandatory failure
of the old restored-$\ii$ sign, paired Fourier reality, all three declared
Bloch transitions, the classical retarded TT step solution, the two
collinear Ward complexes and mutants, the $\phi^4$ obstruction and projected
repair, the nonzero matter Poisson bracket, and the fermion divided differences.
It does not instantiate a full three-dimensional matter complex or execute a
quantum Kubo calculation.  Its successful run reports 55 exact
symbolic/rational checks and four floating-point regressions.  The floating
checks sample generic fibres and random scalar/fermion data; they establish no
universal statement.

The exact cost is fixed-size symbolic algebra plus a four-site matter system;
the numerical cost is constant and negligible.  The certificate is read-only.

\begingroup
\renewcommand{\refname}{References for the matter-coupling proof}
\begin{thebibliography}{9}
\bibitem{toe-r7-matter-Leclerc}
M.~Leclerc, ``Canonical and gravitational stress-energy tensors,''
\href{https://arxiv.org/abs/gr-qc/0510044}{arXiv:gr-qc/0510044} (2006).

\bibitem{toe-r7-matter-BuonannoCellaCurci}
A.~Buonanno, G.~Cella, and G.~Curci, ``Lattice energy-momentum tensor with
Symanzik improved actions,''
\href{https://arxiv.org/abs/hep-lat/9506008}{arXiv:hep-lat/9506008} (1995).

\bibitem{toe-r7-matter-GiustiPepe}
L.~Giusti and M.~Pepe, ``Energy-momentum tensor on the lattice:
non-perturbative renormalization in Yang--Mills theory,''
\href{https://arxiv.org/abs/1503.07042}{arXiv:1503.07042} (2015).
\end{thebibliography}
\endgroup
\endgroup
